\documentclass[11pt,a4paper]{article}
\usepackage[T1]{fontenc}
\usepackage{lmodern}
\usepackage[margin=27mm]{geometry}
\usepackage{amsmath,amssymb,microtype}
\usepackage[hidelinks]{hyperref}
\usepackage{enumitem}
\setlist{itemsep=4pt,topsep=5pt}
\emergencystretch=2em
\clubpenalty=10000
\widowpenalty=10000
\newcommand{\Ostar}{O^{*}}
\newcommand{\draftlabel}[1]{\texttt{\detokenize{#1}}}
\title{Dense Max 2-SAT\\Revision Plan and Exposition Standards}
\author{}
\date{20 October 2026}

\begin{document}
\maketitle

This note records the agreed organization and writing standards for our
dense Max 2-SAT paper. The immediate objective is a readable theoretical
arXiv version centered on the technique and four principal results,
arranged by increasing complexity of the method. Practical extensions
and application studies belong to a later stage directed toward a JAIR
submission. Additional theoretical material remains available in the
working notes and appendices while the main exposition stabilizes.
This is an editorial plan; the
proposals explicitly marked for verification below are not new proved
claims. The standards apply to subsequent revisions as well as the next pass.

\section{Scope and organizing principle}

Organize the paper around the techniques that make accumulated pruning
constraints useful. Whether an input is random is a property of an
application, rather than the main organizing principle. In particular,
deterministic graph families and random graph families can verify similar
structural conditions and should be explained together when they use the
same method.

Keep the principal objective visible: a deterministic exact algorithm with
a fixed exponential saving over $2^n$, using polynomial space, on the
specified families. State which fixed class parameters determine the saving.
Do not confuse this objective with an unrestricted-space upper bound, or
with an improvement whose saving tends to zero as density grows.

The four principal results should concern natural input families. General
distributional or matrix conditions belong in supporting results when they
are needed to prove those family statements. Do not promote every useful
technical lemma or specialization to another headline theorem.

\section{Publication stages and the first release}

\subsection{First objective is a focused theoretical arXiv paper}

Disseminate the technique and selected theoretical results promptly.
The four-part plan and the exposition standards remain unchanged.
Limit the first release to a coherent theoretical contribution with
precise statements, executable algorithm descriptions, complete proofs,
and the related work needed to place the results accurately.

Natural graph families such as Paley graphs, qualifying Ramanujan graphs,
and SBM are theoretical applications of the method. They remain in scope
without requiring an empirical demonstration of practical relevance.
A short motivation and a short concluding outlook are sufficient for
the broader application ambitions.

The first release does not depend on a practical solver implementation,
benchmark experiments, a survey of application domains, new application
classes, or Lean formalization. Optional improvements, including
unification of the sequential and two-round proofs, should not become
prerequisites for dissemination. If such a change remains unresolved,
retain the existing proof route, check it for the selected statement,
and present its dependencies clearly.

The release milestone is a self-contained account of the selected
results: the algorithms use exactly the information stated, the proofs
establish the claimed bounds and polynomial space, and the introduction,
abstract, and result table accurately describe those claims. Reducing
scope should make this milestone attainable sooner.

\subsection{Later objective is a theory and applications paper}

The longer-term aim is a JAIR submission combining the theoretical
foundation with useful practical extensions and input classes arising
in real applications. Identify concrete encodings, determine whether
their instances satisfy the relevant assumptions or suggest useful
relaxations, and evaluate the resulting algorithms where appropriate.
These are later research tasks. They should not appear as established
practical contributions in the first version.

Keep application ideas and solver extensions in working notes so that
they remain available for this stage. Their eventual inclusion should
follow the evidence and the coherence of the extended paper.

\subsection{Material to defer from the current draft}

Use three destinations: the main theoretical account, appendices needed
to support it, and separate working material for later development.
An appendix is not automatically a destination for every existing result.
The following are initial scope decisions, rather than instructions to
delete research material.

\begin{itemize}
\item Defer the extended modeling and practical-application discussion
in \draftlabel{sec:applications}. Preserve it in working notes; retain
only concise motivation in the first paper.
\item Defer region comparisons, caching, and the proposed unified solver
in \draftlabel{sec:cache} and \draftlabel{sec:unified}. They add mechanisms
outside the four selected results.
\item Reduce the extended redundancy and randomized-exploration agenda
to a short outlook. Preserve the full discussion as future research
material.
\item Keep repeated interaction profiles, bicliques, and bounded-block
compression in the working inventory initially. Include an additional
theoretical result only if it materially clarifies the selected technique
without creating another substantial line of exposition.
\item Keep general random models, alternative analyses, and sharper
constants only where they support or briefly clarify a selected result.
They need not all appear in the first arXiv version.
\item Keep the development record, exploratory finite checks, and formal
verification priorities in author working notes. Distinguish such records
from proof details required by the published statements.
\end{itemize}

All proof machinery needed by a selected theorem remains in the paper or
its supporting appendices, with standard external results cited precisely.
In particular, the hidden-community algorithm's spectral preprocessing
and its required computational guarantees remain part of the theoretical
contribution.

\section{The four parts of the main presentation}

\subsection{Pruning constraints and relative literal bias}

Present the common comparison rule, correctness, propagation, simplification,
and summation. Establish the relevant enumeration fact before its first use.
Use the relative literal-bias theorem as the first principal application:
substantial local preferences create useful requirements in the opposite
branches.

This foundation is substantially in place. Preserve it during the next
revision except for necessary corrections and notation alignment. A later
consistency pass should reconcile it with the reorganized sections.

Keep the actual degree and bias assumptions visible. A relative gap is
proportional to a variable's degree; it need not be proportional to $n$.
Do not describe the theorem as requiring only literal gaps if its statement
also requires control of the degree spread.

\subsection{Normalization and symmetric random signs}

The second principal result is a fixed exponential saving for familiar
random clause models on an arbitrary prescribed set of variable pairs.
Lead with one uniformly signed clause per selected pair and independent
inclusion of the four possible clauses with probability $1/4$. Emphasize
the absence of a density or regularity requirement for these models.
Keep the assumption that any additional units are fixed independently of
the binary draw.

Revise the intuition around two separate tasks. Normalization puts
constraints of different sizes on comparable scales before addition.
Random signs prevent the effects of shared variables from remaining
coordinated across all the constraints. Explain these tasks before
presenting degree-dependent estimates.

Use one running example. A larger numerical gap need not mean stronger
pruning when the range of typical fluctuations is also larger. Explain
that the algorithm weights a constraint using the degree in the original
observed graph, rather than estimating its current strength exactly.

An independently sampled support, including an SBM support subsequently
decorated with independent random literal signs, is an immediate
application by conditioning on that support. It does not need another
algorithm or a separate extended discussion.

Place the general complementation-invariant pair law, further sampling
parameters, and sharper constants in an appendix. The proof supporting
the familiar main-text cases must remain complete.

\subsection{Assignment counts and a common systematic interaction}

Explain why some pair interactions have a common systematic contribution,
and how guessing the total number of ones accounts for it. Introduce the
comparison allowance needed for correctness. Then explain the structural
conditions controlling the remaining correlations.

The principal application is Max-Cut on regular graphs satisfying the
stated spectral and density conditions. Present Paley graphs and qualifying
nonbipartite Ramanujan graphs immediately as examples of that result,
retaining their qualifications. They need not receive separate theorems.

Use a supporting deterministic accumulation result to connect this
application with independent random interactions having a common mean.
Max-Cut on $G(n,p)$ is a useful bridge to the next part: there is one
expected pair interaction, and hence one total count to consider.
Include its supporting verification before treating it as an established
corollary of the reorganized argument.

\paragraph{Proposal requiring verification.}
Prefer one fixed-round analysis in the main presentation if it supports
all the selected applications. In particular, investigate using two prefix
rounds for the independent common-mean models as well as regular-graph
Max-Cut. Check the full-matrix bound, the interaction between rounds, the
comparison allowance, and the resulting constants. The existing sequential
proof and the round proof have different hypotheses; do not silently
identify them or promise unchanged numerical bounds.

Do not force the arbitrary-support random-sign theorem through this
matrix framework. Its direct exposure argument handles irregular supports
and has an independent role.

\subsection{Several systematic interactions and hidden communities}

The fourth principal result is Max-Cut on SBM graphs with hidden
communities, under the existing fixed-community and variance assumptions.
The theorem's algorithm receives the graph and the stated fixed class
parameters, without the generating partition or edge probabilities.

First explain the supplied-community case: guessing the number of ones
in each group extends the preceding total-count idea. Use it as an
intermediate explanation. Then explain how approximate spectral
directions and projection guesses replace the unavailable information
in the current hidden-community algorithm. Recovering the community
labels is not required.

Keep the distinction between two uses of SBM explicit. Random signed
clauses on an SBM support belong to the arbitrary-support result.
The fixed complementary clauses encoding Max-Cut on an SBM graph belong
to this fourth part. Supplied and hidden communities are different
information settings for the latter distribution.

\section{A consistent pattern within each part}

Use the following sequence, allowing short definitions immediately before
their first necessary use.

\begin{enumerate}
\item \textbf{Intuition for why the result is possible.}
Identify the useful feature of the input, the obstacle it resolves, and
the source of the search-space saving. Use a concrete example or a path
in the branching tree. Do not begin with the full technical apparatus.

\item \textbf{Definitions needed by the algorithm.}
Introduce only the new notation required here. Refer to the common
framework for established definitions. Keep notation used only by a
later extension in the section that needs it.

\item \textbf{The algorithm.}
Give a short operational description followed by compact but specific
pseudocode. Refer to previously stated procedures and lemmas rather
than repeating their implementation.

\item \textbf{Correctness.}
Explain why an optimum remains represented, including the treatment of
ties, comparison allowances, simplification, and the interaction with
previously added constraints. Separate this argument from the time bound.

\item \textbf{The computational result.}
State a self-contained complexity theorem for the input family. Define
its parameters, assumptions, probability guarantee, and space requirement.
The theorem states what can be computed and at what cost; the algorithm
is the tool used in its proof.

\item \textbf{A roadmap for proving the running-time bound.}
Explain what must be counted and which estimates give the saving.
Distinguish this roadmap from the earlier intuition about why the
phenomenon occurs.

\item \textbf{The main proof.}
Organize it into meaningful steps and show how they establish the stated
bound. Keep the principal argument readable without repeatedly leaving
it for technical derivations.

\item \textbf{Supporting lemmas and applications.}
Place technical estimates after the main proof or in an appendix when
that improves its flow. Introduce every lemma and theorem with a short
explanation of its role, naming the argument that uses it. Facts needed
to define the algorithm must already have been stated.
\end{enumerate}

The section should remain coherent if the reader first follows only the
intuition, algorithm, theorem, and proof roadmap. The full proof must
then supply every dependency needed to justify that account.

\section{Writing and notation standards}

\subsection{Clarity takes priority over brevity}

Longer proofs are welcome when the extra explanation helps the reader.
Use connected prose, meaningful intermediate steps, and consistent
notation. Avoid compressing several ideas into a sentence merely to
shorten a proof. Preserve the author's intended style and message when
editing prose.

Use familiar mathematical language whenever it is sufficient. Say
``sum of Boolean literals'' when that describes the expression. Introduce
general coefficients, affine forms, or matrix notation when they perform
a necessary role, and explain that role at introduction.

Define every symbol, operator, and term before use. In particular, explain
``unassigned variable,'' ``preferred value,'' ``prefix,'' ``suffix,'' and
any order on assignments. Do not use ``free variable'' or ``bad value''
without a definition. Prefer descriptive terms that remain meaningful
from the algorithm's point of view.

Use $C$ for one pruning constraint and $\mathcal C$ for the collection.
Define the summed constraint $C_\Sigma$ in a displayed equation and
refer to that equation consistently. Distinguish coefficients arising
from repeated literal occurrences from weights assigned to entire
constraints. Explain multiplicity as counting appearances.

Write ``by the Cauchy--Schwarz inequality'' and display the particular
inequality being applied. Apply the same principle to other standard
tools: give the reader enough information to check the step.

\subsection{Statements must describe their actual scope}

Define the input model positively and precisely: simple formulas,
bounded multisets, units, or graphs, as appropriate. State assumptions
where they are used rather than relying on an earlier informal
description. Theorem statements should define parameters such as
$\varepsilon,\eta,R,k,\chi$ and explain which are fixed.

State substantive hypotheses in the theorem itself instead of making
the reader reconstruct them from an equation reference. Keep technical
conditions out of headline family descriptions where possible, while
retaining all conditions in the formal result.

Distinguish a linear lower bound from an $O(n)$ upper bound. Distinguish
degree, clause density, and the number of clauses. Define the population
over which an average is taken and whether clauses or distinct neighbors
are counted.

For a random-input result, specify the sampling process, the independent
choices, and what is fixed first. State whether units may depend on the
sample. Separate correctness on every input from an expected or
high-probability running-time guarantee. Do not convert a high-probability
bound into an expected bound without an argument.

Keep proved statements, proposed refinements, and open directions
distinct. Do not claim novelty or completion of an independent proof
audit on the basis of an exposition revision.

\subsection{Intuition should explain a mechanism}

A statement such as ``typical assignments cannot compensate for all
requirements'' needs an explanation: bounded influence, opposing signs,
a structural correlation bound, or another identified mechanism. Use a
small example to make the mechanism visible.

When using a uniform random path, say that this is an analysis of a
deterministic search tree. A literal contributes zero or one; one half is
its expectation under the specified uniform assignment. Do not confuse
randomness of the path with randomness of the input.

Do not assume that generated constraints are independent. Do not imply
that fresh variables or disjoint neighborhoods are necessary when the
argument actually controls the signs or total influence of overlaps.
Explain how the proof handles exceptional paths and sums work over all
prefixes.

Introduce normalization before repeated comparisons involving
$\sqrt d$. Explain that positive scaling preserves an individual
constraint but changes its contribution to the sum. Treat normalization
and removal of a systematic interaction as separate operations.
Avoid calling useful pruning a ``disadvantage''; describe the requirement
imposed on the branch or the resulting saving explicitly.

\section{Algorithm presentation standards}

Pseudocode should be high-level enough to follow and precise enough to
execute in principle. Use the \texttt{algorithm} and \texttt{algorithmic}
packages. Comments should be short explanations in a smaller font,
introduced by \texttt{//}; the current convention is
\begin{verbatim}
\renewcommand{\algorithmiccomment}[1]{{\footnotesize\itshape // #1}}
\end{verbatim}

\begin{enumerate}
\item \textbf{Follow execution order in the description.}
The reader should encounter the substantive steps in the order in which
the algorithm performs them. Introduce the input, fixed variable order,
preprocessing, and other setup before describing the search. Explain
rejection and stopping conditions before the operations they may prevent.
For nested loops, introduce the outer guess and current assignment before
explaining computations inside them. Define each quantity before its
first use, and align the prose with the pseudocode. Omit routine details
that an algorithms reader can infer.

\item \textbf{Identify the input and all state.}
State parameters explicitly or define a shared-data block before the
procedures. Distinguish fixed data, shared updates, and values local to a
recursive call. Explain $T$, $M$, and $\alpha^\star$, where they are
initialized, and which calls can change them. Define score offsets and
how the returned assignment and score include them.

\item \textbf{Specify branching choices concretely.}
State the variable order and how the preferred value is computed.
``Increasing labels'' means the original input indices if that is the
intended order. Distinguish original graph degrees from residual degrees.
If a theorem permits any fixed order independent of the sample, say so;
do not infer adaptive-order guarantees without the appropriate proof.
In round algorithms, state which earlier assignments the preferences use.

\item \textbf{Separate initial preprocessing from node operations.}
Use \textsc{Preprocess} for the initial pair reductions and
\textsc{Simplify} for the justified local rules and constraint propagation
at search nodes. Do not repeat a reduction lemma already provided by
the common framework. Apply these procedures only as permitted by the
particular algorithm and its analysis: adding forced assignments to a
fixed-prefix random-search analysis requires justification.

\item \textbf{Respect constraints and comparison scope.}
Explain how simplification and learned constraints jointly preserve an
optimum. Individual correctness of two rules does not by itself establish
their compatibility. Branch exploration order is separate from the
mathematical preference and tie convention; do not impose a first-branch
requirement when the correctness invariant does not need one.

\item \textbf{Use only information supplied or computed.}
For the random-formula algorithm, construct the observed support from
the input formula before preprocessing. The latent generating graph is
not an input. For SBM, distinguish supplied structural data from hidden
data and account for discovering the information actually used.

\item \textbf{Explain what a test guarantees.}
Failure of the summed-constraint test certifies infeasibility, but passing
it need not imply joint feasibility. Likewise, satisfying $C_\Sigma$
need not satisfy every member of $\mathcal C$. Explain why enumerating
these extra candidates preserves correctness and how their number is
bounded.
\end{enumerate}

Do not introduce caching as an implicit part of the core algorithm.
Make the polynomial-space implementation explicit: depth-first traversal
and sequential processing of guesses and candidates, rather than storing
the whole search space.

\section{Proof and linear algebra standards}

Separate preservation of an optimum, the count of surviving assignments,
enumeration time, arithmetic cost, and space. A small number of solutions
to an arbitrary constraint system does not by itself yield an efficient
enumerator. Invoke the proved enumeration property of the particular
summed inequality when that is the argument.

For probabilistic proofs, specify what is exposed or conditioned on,
why the remaining random choices have the required distribution, and
which dependencies remain. When introducing a martingale, define the
available information, verify conditional mean-zero increments, and
justify the increment or variance bounds used by the concentration
inequality. Explain transitions from probability over assignments to
counts, and from joint input--assignment estimates to guarantees for
most inputs.

Use supporting lemmas for technical estimates that would interrupt the
main proof. Keep comments on stronger variants, adaptive choices, or
possible extensions in a separate discussion unless they are needed
for the stated conclusion.

Introduce linear algebra when its role becomes necessary to the
presented argument. Keep the following distinctions explicit:
\begin{itemize}
\item Symmetric random signs on arbitrary supports use scalar moment
and concentration arguments in the current proof.
\item Count-adjusted independent interactions have a combinatorial
algorithm but currently use matrix bounds in the analysis.
\item Regular-graph Max-Cut uses its spectral hypothesis in the proof;
the algorithm does not calculate eigenvectors.
\item Supplied-community SBM uses counts and gains algorithmically,
with matrix estimates in the proof.
\item Hidden-community SBM currently performs spectral preprocessing
as part of the algorithm.
\end{itemize}
Moving a proof to an appendix does not remove its mathematical
dependencies. A proposed elementary replacement must be proved before
claiming that the relevant technique has been eliminated.

\section{Appendices and revision sequence}

The first arXiv version's appendices primarily contain the complete
supporting proofs for the four principal results. Include a general
distributional statement, parameter refinement, or alternative analysis
when it supports those results or gives a short useful clarification.
The broader inventory of proved results and exploratory ideas remains
available in separate working material; it need not all accompany the
submission. Keep proposed extensions visibly separate from established
results wherever they are retained.

The next passes should proceed as follows.
\begin{enumerate}
\item Revise Part 2 around normalization and the effect of random signs,
using the agreed section pattern and the familiar sampling models.
\item Rebuild Part 3 around counts and the selected deterministic
Max-Cut result. Verify any change of analysis before replacing an
existing proof; place additional random interaction results appropriately.
\item Rebuild Part 4 from supplied community counts to the actual
hidden-community algorithm, retaining its assumptions and computational
steps.
\item Align the introduction, abstract, result table, transitions, and
cross-references with the four-result scope. Remove superseded family
descriptions and separate the deferred application and research material.
Preserve author-written prose through careful, local edits.
\item Check every retained statement against its proof, algorithm,
input information, and space bound; compile and inspect the manuscript.
Complete the focused arXiv version before expanding the scope toward
practical extensions and the later journal submission.
\end{enumerate}

The existing Lean verification plan remains a separate resource. Reconcile
its theorem priorities with the stable manuscript before beginning
formalization, as previously agreed.

\subsection{Map to the current draft}

The source is \nolinkurl{dense_max2sat_paper.tex}. Use labels rather than
theorem numbers while reorganizing:
\begin{description}[style=nextline]
\item[Part 1]
\draftlabel{sec:framework}, \draftlabel{sec:bias},
\draftlabel{lem:pruning}, \draftlabel{lem:constraint_sum}.
\item[Part 2]
\draftlabel{sec:random}, \draftlabel{thm:weighted},
\draftlabel{cor:support}, \draftlabel{lem:observed_pairs},
\draftlabel{lem:enumeration}, \draftlabel{eq:weighted}.
\item[Part 3]
\draftlabel{sec:countalgorithm}, \draftlabel{cor:maxcut},
\draftlabel{thm:randomclauses},\\
\draftlabel{res:counts},
\draftlabel{prop:pairmembership}.
\item[Part 4]
\draftlabel{sec:sbm}, \draftlabel{thm:sbm},
\draftlabel{cor:planted}, \draftlabel{app:discovery}.
\end{description}
This map identifies source material for the revision; it does not freeze
the final labels or require the present section boundaries to survive.

\end{document}
